\subsection{A specific Borcherds theta construction is the cubic discriminant}

The automorphic ring in P-F37 has a known theta construction for a
particular nonzero invariant. This is substantially more specific than
asking for an unidentified theta realization of an entire multiplicity
space. The construction and its divisor are known; the following
corollary records their exact compatibility with the plethysm grading.
No historical priority is claimed for the construction or this
identification.

\paragraph{Known geometric input.}
Let $\mathbb B$ be the cubic-threefold ball of dimension ten, with
nodal and chordal arrangements $\mathcal H_\Delta$ and $\mathcal H_c$.
Casalaina-Martin--Laza, Proposition~7.5, printed p.~25,
construct a weight-$48$ form $\phi$ from the Borcherds form
$\Phi_{12}$ by an $E_6$ quasi-pullback and restriction to this ball.
Its orders along these two arrangements are respectively $3$ and $84$.
Their Lemma~7.7 and Corollary~7.8, p.~26, give ramification
orders $3,6$ and the quotient divisor $H_\Delta+14H_c$.
The source is \url{https://arxiv.org/pdf/0710.5329}.

Write $U=\mathbb C^5$, $\mathscr S=\operatorname{Sym}^3U^*$,
and let $\Delta_{\rm cub}$ be the usual discriminant polynomial
on $\mathscr S$, with an arbitrary fixed nonzero normalization.
The period-map isomorphism of P-F37 will be denoted by
\[
 \mathcal P:\bigoplus_k\mathcal M_k
       \longrightarrow \mathbb C[\mathscr S]^{SL(U)};
 \qquad\deg\mathcal P(g)=\frac53\operatorname{wt}(g).
 \tag{P-F41a}
\]

\begin{corollary}[P-F41: an identified theta image in the actual cubic plethysm algebra]
\label{cor:thetadiscriminant}
The form $\phi$ has trivial multiplier on the group used in
(P-F41a), and
\[
             \mathcal P(\phi)=c\,\Delta_{\rm cub}
             \quad\text{for some }c\in\mathbb C^*.
 \tag{P-F41b}
\]
Consequently this particular multiplicative Borcherds theta
construction identifies a nonzero vector in the actual graded piece
\[
 T^{(5)}_{16}
    =\left(\operatorname{Sym}^{80}
                    (\operatorname{Sym}^3\mathbb C^5)\right)^{SL_5},
 \tag{P-F41c}
\]
with automorphic weight $48$, ordinary coefficient degree $80$,
and stretching degree $16$. Under the full $GL_5$ action its
character is $\det^{48}$. It identifies the polynomial subalgebra
$\mathbb C[\phi]$ with $\mathbb C[\Delta_{\rm cub}]$,
including their products and these gradings.

Here ``Borcherds theta construction'' has a precise meaning:
$\Phi_{12}$ is the multiplicative lift associated with the
regularized orthogonal theta integral for the input
$\delta(\tau)^{-1}=q^{-1}+24+\cdots$ on $II_{2,26}$,
where $\delta$ is the elliptic modular discriminant.
The integral gives the logarithm of its Petersson norm;
$\Phi_{12}$ is not asserted to be the value of an ordinary
linear convergent theta integral. The subsequent operations
are quasi-pullback and restriction.
\end{corollary}

\begin{proof}
We first check the multiplier instead of discarding it implicitly.
Let $\Lambda$ be the Eisenstein primitive cohomology lattice of
a cyclic cubic fourfold and $\Gamma_{\rm mon}$ its monodromy group.
Allcock--Carlson--Toledo, Lemma~2.3, printed p.~11, sends a
nodal meridian to an exact order-three complex reflection fixing
its mirror pointwise. Their Theorem~8.1 and proof, p.~74,
identify $\Gamma_{\rm mon}$ with the subgroup of
$\operatorname{Aut}_E(\Lambda)$ having determinant in $\mu_3$;
the proof also records that these meridians generate the
monodromy image. These are primary statements in
\url{https://arxiv.org/pdf/math/0608287}.

Suppose initially that $\phi$ has a scalar multiplier $\chi$.
This allows for the character in a Borcherds automorphy statement.
View it as a homogeneous function of degree $-48$ on the
negative cone over the ball. For a nodal reflection $r$, choose
a generic point of its mirror and a linear slice of that cone
defined by a nonzero linear functional invariant under $r$.
Such a slice exists because $r$ fixes the mirror pointwise.
On this slice choose normal coordinate $z$ with
$r(z)=\omega z$, where $\omega^3=1$ and $\omega\ne1$.
Its order-three zero gives a local expansion
\[
               \phi(z,y)=z^3 u(z,y),\qquad u(0,y_0)\ne0.
\]
The automorphy equation $\phi(rZ)=\chi(r)\phi(Z)$,
compared in its first nonzero normal coefficient at $y_0$,
gives $\chi(r)=\omega^3=1$. It follows that
$\chi$ is trivial on $\Gamma_{\rm mon}$.
If the automorphy group is taken to be the full unitary group,
the same conclusion holds there: a unitary lattice determinant
is an Eisenstein unit, hence belongs to $\mu_6$; the subgroup
just described has index two and the other coset is represented
by $-I$, because $\operatorname{rank}_E\Lambda=11$.
Homogeneity gives $\phi(-Z)=(-1)^{-48}\phi(Z)=\phi(Z)$.

For completeness, this also matches precisely the group in
Looijenga--Swierstra rather than merely its projectivization.
Their group consists of the $\mu_3$-equivariant isometries
of full middle cohomology fixing the polarization class
(Section~3, printed p.~10). It contains the monodromy group.
The proof of Allcock--Carlson--Toledo Theorem~8.1 explicitly
notes that $-I$ on primitive cohomology does not extend to
full cohomology while fixing that class. Since the full
unitary group has just the two cosets above, its extendable
subgroup is exactly $\Gamma_{\rm mon}$. Thus the form is
invariant for the group of (P-F41a), with no need to take a power.

Apply the actual graded ring isomorphism of
Looijenga--Swierstra Theorem~3.1, p.~12, and its grading
calculation on pp.~10--11, as in P-F37. The proof on p.~13
identifies the automorphic and invariant line bundles over
the stable locus, so the identification also preserves
vanishing at a stable nodal cubic.
It sends the nonzero form $\phi$ to a nonzero invariant
$I\in\mathbb C[\mathscr S]^{SL_5}$ of degree
$(5/3)48=80$. A holomorphic form on the ball is in particular
allowed in the arrangement-meromorphic ring of that theorem.
At a nodal cubic, the extended period map meets the nodal
mirror; this discriminant identification is also explicit
in Allcock--Carlson--Toledo Theorem~1.1, printed p.~3.
Since $\phi$ vanishes there, its pullback $I$
vanishes on the nodal discriminant locus.

We specify why this locus is dense in the whole discriminant.
The incidence variety
\[
 \mathscr I=\{(F,p)\in\mathscr S\times\mathbb P^4:
                         dF(p)=0\}
\]
is a rank-$30$ vector bundle over $\mathbb P^4$: after sending
$p$ to $[1:0:0:0:0]$, the five conditions set the coefficients
of $x_0^3,x_0^2x_1,\ldots,x_0^2x_4$ to zero.
It is irreducible of dimension $34$ and projects properly
onto the discriminant. Requiring singularities at two
distinct points imposes ten independent conditions. Indeed
move the ordered pair to $[1:0:\cdots:0]$ and
$[0:1:0:\cdots:0]$; the two lists of five cubic monomials
are disjoint. The ordered-pair incidence therefore has
dimension $8+(35-10)=33$.
The locus in $\mathscr I$ where the quadratic part at $p$
is degenerate is also proper closed. To see properness,
\[
        F_0=x_0\sum_{i=1}^4x_i^2+\sum_{i=1}^4x_i^3
 \tag{P-F41d}
\]
has its only singular point at $[1:0:\cdots:0]$, where its
affine quadratic part is nondegenerate. In fact a critical
point with $x_0=0$ forces all $x_i=0$. Otherwise set
$x_0=1$; every $x_i$ is $0$ or $-2/3$, and
$\partial F_0/\partial x_0=\sum_i x_i^2=0$ forces
all four to vanish. Thus the exceptional loci just described
have images, or image closures, of dimension at most $33$.
Away from them a singular cubic has a unique ordinary node.
The nodal locus is therefore dense in the dimension-$34$
discriminant hypersurface.

The exact algebraic discriminant facts are
Bus\'e--Jouanolou, Proposition~4.7(i), printed p.~40,
and Theorem~4.26, p.~53, in
\url{https://arxiv.org/pdf/1210.4697}.
Over $\mathbb C$ their generic discriminant is a prime
polynomial and its coefficient degree is
$5(3-1)^4=80$. It cuts out the singular-hypersurface locus.
Consequently $I$, vanishing on the dense nodal locus,
vanishes on its whole irreducible closure and is divisible
by $\Delta_{\rm cub}$. Both polynomials have degree $80$,
so $I=c\Delta_{\rm cub}$ with $c\ne0$.
This proves (P-F41b). The change-of-variables formula in
the same source, Proposition~4.13, pp.~42--43, gives
$\Delta_{\rm cub}(F\circ g)=\det(g)^{3\cdot2^4}
\Delta_{\rm cub}(F)$, proving the stated $GL_5$ character.
The regrading by five is the actual Cartan/plethysm ring
identification already made in P-F19 and P-F37, so this
proves (P-F41c), not merely an equality of dimensions.
Since the discriminant is a nonconstant polynomial,
the homomorphism $\mathbb C[t]\to\mathbb C[\Delta_{\rm cub}]$
sending $t$ to $\Delta_{\rm cub}$ is injective:
if $p(\Delta_{\rm cub})=0$ for a nonzero $p\in\mathbb C[t]$,
factor $p$ over $\mathbb C$ and use that the polynomial
ring is a domain to obtain the contradiction that
$\Delta_{\rm cub}$ is constant.
The ring isomorphism then gives the assertion about
$\mathbb C[\phi]$ and all its products.

We finish by verifying the precise meaning of the theta
construction. For $M=II_{2,26}$ the weakly holomorphic
input $\delta^{-1}$ has the required weight $-12=1-26/2$
and has integral Fourier coefficients at nonpositive
exponents. Borcherds, \emph{Automorphic forms with
singularities on Grassmannians}, Section~6, printed p.~22,
defines the theta integral by truncation and meromorphic
regularization, taking a constant term at the regularizing
parameter zero. Theorem~13.3, pp.~48--49, gives its
multiplicative lift of weight $24/2=12$; item~(3)
relates the logarithm of the lift's norm to that integral.
The input's only negative term is $q^{-1}$, giving the
order-one root divisors of $\Phi_{12}$.
The identification of this product with $\Phi_{12}$ is
explicit in Gritsenko, \emph{24 faces of the Borcherds
modular form $\Phi_{12}$}, Section~1, equation~(1),
p.~4, and Section~3.2, pp.~11--12.

The primitive cubic-fourfold lattice embeds into
$II_{26,2}$ with orthogonal complement $E_6$.
The operation in Casalaina-Martin--Laza Proposition~7.5
first removes the $36$ root-hyperplane factors containing
the smaller domain and then restricts to it; restriction
to the cubic-threefold ball follows. Removing a linear
factor increases automorphic weight by one, so the resulting
weight is $12+36=48$. These are exactly the geometric
operations producing the $\phi$ used above. Thus the
theta origin and the invariant identification refer to
the same form. This completes the proof.
\end{proof}

\paragraph{Exact remaining comparison with the supplied method.}
The supplied \emph{Remark}, printed p.~3, already mentions
$O(2,n)$, so the preceding orthogonal construction is not
excluded merely by its group name. Its item~8 on p.~4
specifies iterated theta lifts between totally definite
unitary groups; items~9--10 and~12 concern their norms,
divisibility and periods. The concrete construction above
instead uses a weakly holomorphic scalar input on $SL_2$,
a regularized multiplicative lift to $O(2,26)$, an $E_6$
quasi-pullback, and restriction to $U(10,1)$.
No identification of these operations with the supplied
iterated lift or its determinant has been established.

The next comparison can now be stated for a named output:
it must recover $\phi$ (equivalently $\Delta_{\rm cub}$),
or an identified larger graded collection containing it,
with the three gradings in (P-F41c) respected.
The existing construction supplies neither a basis for
every plethysm multiplicity space nor a uniform positive
counting formula, efficient oracle, or PH2/SH estimate.
It concerns the fixed inner representation
$\operatorname{Sym}^3\mathbb C^5$.

The regularized-theta primary sources are
\url{https://arxiv.org/pdf/alg-geom/9609022}, Section~6 p.~22,
Theorem~13.3 pp.~48--49, and
\url{https://arxiv.org/pdf/1203.6503}, equation~(1) p.~4
and Section~3.2 pp.~11--12.
